\section{Conclusion}
\label{sec:conclusion}

Probe takes the shape of a researcher's afternoon and compresses it. A rough premise goes in; three divergent research programs come out; two of them get blocked with specific evidence-anchored findings; one survives a three-reviewer panel that refuses to collapse its disagreements; a provenance-tagged guidebook emerges with a linter running on it before it is allowed to leave the run directory. On the demo premise, the pipeline surfaced an announcement-duration confound neither the authors nor the shallow audit had caught, and a Managed Agents deep-audit session ran twenty-one tool calls to measure claims the single-prompt audit could only reason about.

The project's load-bearing commitment is the separation of rehearsal from evidence. Every paragraph that could be mistaken for a finding carries a tag that says otherwise; the linter rejects the build when a \texttt{[SIMULATION\_REHEARSAL]}-tagged paragraph uses evidence language; a \texttt{[HUMAN\_REQUIRED]} block at the end of every guidebook forces the handoff to be explicit. The commitment is structural, not ornamental: a run that violates it fails the lint and does not ship.

What Probe does not claim to be is an autonomous scientist. It refuses the framing. A researcher who uses the tool is working with a skeptical colleague who can produce a passable first draft of a study design in half an hour. That draft is not the study. The study still needs human judgment about which reviewer's objection to prioritize, whether a blocking pattern fired correctly or mis-matched a manipulation-under-test, and whether the simulated walkthrough's friction points generalize to real participants.

The ablation result the paper reports against its own initial claim is the illustrative moment. We had expected Opus 4.7 to uniquely catch the confound that the Sonnet 4.6 comparison revealed both models could find. The honest write-up tempered the capability claim: Opus provides more precise evidence quantification at higher cost; Sonnet is substantively capable on the same prompt. This is the kind of correction Probe was built to make on other people's studies, and the paper's authors are not exempt.

Future work falls along three lines. First, a multi-trial hallucination evaluation with blind scoring to convert the current single-trial pass into a capability claim the paper can defend. Second, a citation-misuse detector that verifies \texttt{[SOURCE\_CARD:id]} paragraphs against the source card's \texttt{allowed\_claims} list --- the one failure mode the provenance linter did not catch on the demo run. Third, a pattern-applicability check that prevents capture-risk patterns from mis-firing against manipulations-under-test --- the cross-domain transfer issue the code-review benchmark exposed.

The project lives at \url{https://github.com/Apolotary/probe-researcher}. The demo run, the ablation report, and the hallucination test output are all checked in under the \texttt{runs/demo\_run/} and \texttt{benchmarks/} directories. The linter, the schemas, and the 16-pattern capture-risk library are in the repository and are small enough to read end-to-end. The defense the project can offer, if a reviewer reads no further than the linter code, is that the infrastructure is small, legible, and does what the paper says it does.
